\section{Introduction}
\label{sec:introduction}

 Due to the advent of deep reinforcement learning (RL) methods that allow the study of many agents in rich environments, multi-agent RL has flourished in recent years. However, most of this recent work considers fully cooperative settings \citep{omidshafiei2017deep,foerster2017counterfactual,foerster2017stabilising} and emergent communication in particular \citep{das2017learning,mordatch2017emergence,lazaridou2016multi,foerster2016learning,sukhbaatar2016learning}. Considering future applications of multi-agent RL, such as self-driving cars, it is obvious that many of these will be only partially cooperative and contain elements of competition and conflict.


The human ability to maintain cooperation in a variety of complex social settings has been vital for the success of human societies.  Emergent reciprocity has been observed even in strongly adversarial settings such as wars \citep{axelrod2006evolution}, making it a quintessential and robust feature of human life.

In the future, artificial learning agents are likely to take an active part in human society, interacting both with other learning agents and humans in complex partially competitive settings. Failing to develop learning algorithms that lead to emergent reciprocity in these artificial agents would lead to disastrous outcomes.

How reciprocity can emerge among a group of learning, self-interested, reward maximizing RL agents is thus a question both of theoretical interest and of practical importance. Game theory has a long history of studying the learning outcomes in games that contain cooperative and competitive elements. In particular, the tension between cooperation and defection is commonly studied in the iterated prisoners' dilemma. In this game, selfish interests can lead to an outcome that is overall worse for all participants, while cooperation maximizes social welfare, one measure of which is the sum of rewards for all agents. 


Interestingly, in the simple setting of an infinitely repeated prisoners' dilemma with discounting, randomly initialised RL agents pursuing independent gradient descent on the exact value function learn to defect with high probability. This shows that current state-of-the-art learning methods in deep multi-agent RL can lead to agents that fail to cooperate reliably even in simple social settings. 
One well-known shortcoming is that they fail to consider the learning process of the other agents and simply treat the other agent as a \emph{static} part of the environment~\citep{hernandez2017survey}. 

As a step towards reasoning over the learning behaviour of other agents in social settings, we propose \emph{Learning with Opponent-Learn-ing Awareness} (LOLA). The LOLA learning rule includes an additional term that accounts for the impact of one agent's policy on the learning step of the other agents. For convenience we use the word `opponents' to describe the other agents, even though the method is not limited to zero-sum games and can be applied in the general-sum setting. We show that this additional term, when applied by both agents, leads to emergent reciprocity and cooperation in the iterated prisoners' dilemma (IPD). Experimentally we also show that in the IPD, each agent is incentivised to switch from naive learning to LOLA, while there are no additional gains in attempting to exploit LOLA with higher order gradient terms. This suggests that within the space of local, gradient-based learning rules both agents using LOLA is a stable equilibrium. This is further supported by the good performance of the LOLA agent in a round-robin tournament, where it successfully manages to shape the learning of a number of multi-agent learning algorithms from literature. This leads to the overall highest average return on the IPD and good performance on Iterated Matching Pennies (IMP). 

We also present a version of LOLA adopted to the deep RL setting using likelihood ratio policy gradients, making LOLA scalable to settings with high dimensional input and parameter spaces. 

We evaluate the policy gradient version of LOLA on the IPD and iterated matching pennies (IMP), a simplified version of rock-paper-scissors. We show that LOLA leads to cooperation with high social welfare, while independent policy gradients, a standard multi-agent RL approach, does not. The policy gradient finding is consistent with prior work, e.g., \citet{sandholm1996multiagent}. 
We also extend LOLA to settings where the opponent policy is unknown and needs to be inferred from observations of the opponent's behaviour. 

Finally, we apply LOLA with and without opponent modelling to a grid-world task with an embedded underlying social dilemma. This task has temporally extended actions and therefore requires high dimensional recurrent policies  for agents to learn to reciprocate. Again, cooperation emerges in this task when using LOLA, even when the opponent's policy is unknown and needs to be estimated. 